% SCAFFOLD, drafted 18 Sep 2026 from the day-one runs. Every number expands from
% tables/numbers.tex, generated by analysis/paper_numbers.py from runs/. Every unmeasured
% result is a visible \blocked{} marker with the target sentence shape; none is a placeholder
% number. Per the author's own standard, the prose is to be retyped in his voice before
% submission. Title is provisional: it states what one model shows and must be re-checked
% against the other models before it is kept.
\documentclass[11pt]{article}
\usepackage[margin=1in]{geometry}
\usepackage{newtxtext,newtxmath}
\usepackage{booktabs}
\usepackage{array}
\usepackage{amsmath}
\usepackage{graphicx}
\usepackage{xcolor}
\usepackage{url}
\usepackage[hidelinks]{hyperref}
\usepackage{natbib}
\newcommand{\qtN}{400}
\newcommand{\qtSizeFp}{3.8}
\newcommand{\qtAccEight}{91.75}
\newcommand{\qtSizeEight}{2.2}
\newcommand{\qtAccSix}{93.50}
\newcommand{\qtSizeSix}{1.7}
\newcommand{\qtDeltaSix}{+1.75}
\newcommand{\qtLostSix}{6}
\newcommand{\qtGainedSix}{13}
\newcommand{\qtPSix}{0.167}
\newcommand{\qtAccFive}{90.75}
\newcommand{\qtSizeFive}{1.5}
\newcommand{\qtDeltaFive}{-1.00}
\newcommand{\qtLostFive}{10}
\newcommand{\qtGainedFive}{6}
\newcommand{\qtPFive}{0.454}
\newcommand{\qtAccFour}{92.00}
\newcommand{\qtSizeFour}{1.3}
\newcommand{\qtDeltaFour}{+0.25}
\newcommand{\qtLostFour}{13}
\newcommand{\qtGainedFour}{14}
\newcommand{\qtPFour}{1.000}
\newcommand{\qtAccThree}{87.00}
\newcommand{\qtSizeThree}{1.1}
\newcommand{\qtDeltaThree}{-4.75}
\newcommand{\qtLostThree}{30}
\newcommand{\qtGainedThree}{11}
\newcommand{\qtPThree}{0.004}
\newcommand{\qtSizeRatioFour}{59}
\newcommand{\qtCompressFour}{34}
\newcommand{\qtRepRunsEight}{3}
\newcommand{\qtRepMeanEight}{91.58}
\newcommand{\qtRepRangeEight}{0.25}
\newcommand{\qtRepListEight}{91.75, 91.50, 91.50}
\newcommand{\qtRepFlipsMaxEight}{1}
\newcommand{\qtRepRunsFour}{3}
\newcommand{\qtRepMeanFour}{91.83}
\newcommand{\qtRepRangeFour}{0.50}
\newcommand{\qtRepListFour}{92.00, 91.50, 92.00}
\newcommand{\qtRepFlipsMaxFour}{4}
\newcommand{\qtMaxRepRange}{0.50}
\newcommand{\qtCliffOverNoise}{10}
\newcommand{\qtRungsCount}{5}
\newcommand{\qtRunsTotal}{9}
\newcommand{\qtMacTg}{35}
\newcommand{\qtMacPp}{650}
\newcommand{\qtGcpTg}{22.0}
\newcommand{\qtGcpPp}{109.6}

\newcommand{\blocked}[1]{\textcolor{red}{\textbf{[BLOCKED, pending measurement: }#1\textbf{]}}}
\newcommand{\bfcl}{BFCL}
\newcommand{\qm}{Qwen3-1.7B}

\title{Flat to Four Bits, Off a Cliff at Three:\\Where Quantization Breaks Agentic Tool Calling}
\author{Akshay Kumar\\Independent Researcher\\\texttt{akumar8@mt.iitr.ac.in}}
\date{}

\begin{document}
\maketitle

\begin{abstract}
\blocked{Written last. Shape: what practitioners run (k-quants), what has been measured for them (prose, knowledge, reasoning), what has not (tool calling); the ladder design; the one-model result with the cliff and the reproducibility figure; the free-form comparison and the verdict on H1; the multi-model result; the deployment table as the deliverable.}
\end{abstract}

\section{Introduction}

Language models reach laptops by being compressed. The weights are stored in fewer bits than they were trained in, and the file shrinks accordingly. The full-precision file for \qm{} is \qtSizeFp{}\,GB; the version most people would actually run, Q4\_K\_M, is \qtSizeFour{}\,GB, about \qtCompressFour{}\% of it. That compression is post-training quantization, and the format nearly everyone uses on a laptop is the family of llama.cpp k-quants \citep{llamacpp}, because that is what Ollama and every desktop front end serve.

The cost of that compression has been measured with care, but on the wrong tasks for one growing use. \citet{kurt2026quant} quantizes Llama-3.1-8B-Instruct into thirteen GGUF formats and scores each on GSM8K, HellaSwag, IFEval, MMLU, TruthfulQA and perplexity. \citet{lotfi2026overthink} measure five reasoning models under GPTQ, AWQ and FlatQuant on mathematics, code and science questions and find accuracy falling as chain-of-thought inflates. Every task in both papers is free-form or multiple choice. None of them is what an agent does.

An agent calls tools. It emits a function name and a set of typed arguments that a program will execute. That output either parses against the schema or it does not; there is no partial credit and no reader to forgive an approximate answer. The intuition that motivated this study is that a function signature has none of the slack that prose has, so compression should break tool calling at a higher bit-width than it breaks free-form accuracy. We pre-registered that as our first hypothesis and set out to find where the floor sits.

We build a controlled ladder, five k-quant formats produced from one FP16 source, and score each rung on the Berkeley Function Calling Leaderboard \citep{bfcl2025}, which compares a model's call to a reference by deterministic abstract-syntax-tree matching. No model judges another model anywhere in the pipeline. We then run the same weights on a free-form control so that the two floors can be compared on identical parameters.

For \qm{} the answer is not the one we expected. Tool-calling accuracy is flat from Q8\_0 through Q4\_K\_M and falls by \qtDeltaThree{} points at Q3\_K\_M, roughly \qtCliffOverNoise{} times the run-to-run noise of the instrument. \blocked{One sentence on where the free-form floor sits on the same weights and whether H1 holds.} \blocked{One sentence on whether the pattern holds across model sizes.}

\paragraph{Contributions.} A controlled k-quant ladder for agentic tool calling, built from one source per model so that only the weights change between rungs (\S\ref{sec:method}). A reproducibility measurement that most quantization papers omit: the same weights run three times, which bounds what a ladder difference can mean (\S\ref{sec:repro}). The floor for \qm{}, located between Q4\_K\_M and Q3\_K\_M by paired testing (\S\ref{sec:ladder}). \blocked{The free-form comparison on identical weights and the resulting verdict on H1 (\S\ref{sec:freeform}).} \blocked{The same measurement across model families and sizes (\S\ref{sec:models}).} A deployment table that answers the question practitioners currently guess at.

\section{Related work}

\paragraph{Quantization evaluated on free-form tasks.} \citet{kurt2026quant} is the closest work and the one this paper extends. One model, thirteen llama.cpp formats from Q3\_K\_S to Q8\_0 against FP16, six free-form or multiple-choice benchmarks, plus throughput and size. A full-text read finds no tool-calling, function-calling, agentic, structured-output or schema-compliance evaluation, and the conclusion asks for exactly what we add: additional model families and sizes under the same protocol. \citet{lotfi2026overthink} take a different axis. Under GPTQ \citep{gptq2023}, AWQ \citep{awq2024} and FlatQuant, five reasoning models lose accuracy while emitting longer chains of thought, and in up to 52\% of failures the correct answer appears mid-trace without being committed to. Two details bear on us. Their compression methods are research methods, not the k-quants that ship. And their failure categorisation, including the overthinking class, uses GPT-5 as a judge; ours has no judge at all.

\paragraph{Tool-calling evaluation.} \bfcl{} \citep{bfcl2025} scores a model's function call against a reference by abstract-syntax-tree matching over the function name, argument names, types and values, across single, multiple and parallel calls, an irrelevance category where no call is correct, and multi-turn trajectories. Because the matcher is deterministic, the same generation always receives the same verdict, which is what lets a ladder difference be attributed to the weights rather than to the scorer.

\paragraph{The gap.} Quantization has been scored on prose, knowledge, instruction following and reasoning, on research compression methods and on the deployed one. It has not been scored on tool calling under either. That is the measurement here.

\section{Method}
\label{sec:method}

\paragraph{Pre-registration.} Hypotheses and thresholds were fixed on 18 September 2026 before the first grid run and are in \texttt{paper/preregistration.md}. Two full-text summaries were appended to that file during the first rung as documentation; no hypothesis or threshold changed after runs began.

\paragraph{The ladder.} Every rung is produced from a single FP16 GGUF of \qm{} \citep{qwen3} with \texttt{llama-quantize}, so the rungs differ in weights and in nothing else. The formats are a subset of the ladder in \citet{kurt2026quant}, chosen to span the range at one rung per bit-width: Q8\_0, Q6\_K, Q5\_K\_M, Q4\_K\_M and Q3\_K\_M. Ollama's library publishes only three of these for this model, which is why the ladder is built rather than pulled.

\paragraph{Serving.} Each rung is served by \texttt{llama-server} with four parallel slots, an 8{,}192-token context and full Metal offload, under a fixed model alias so that the harness sends identical requests to every rung. Generation is at temperature 0.001. Measured throughput is about \qtMacTg{} tokens per second per slot and \qtMacPp{} tokens per second of prompt evaluation on an Apple M1 Max.

\paragraph{Scoring.} \bfcl{} v4, \texttt{simple\_python} category, \qtN{} questions, deterministic AST matching. Differences between rungs are tested pairwise on the same \qtN{} questions with McNemar's exact test \citep{mcnemar1947}, which uses only the discordant pairs and is the appropriate test when two conditions are scored on identical items.

\paragraph{Reproducibility control.} Two rungs, Q8\_0 and Q4\_K\_M, were run three times each with identical weights and settings. This bounds the run-to-run variation of the instrument, without which a ladder difference cannot be interpreted.

\paragraph{Free-form control.} \blocked{The same five GGUFs on 400 GSM8K test questions drawn at seed 42, full generations logged, scored under a strict extractor and a lenient one. Generation is running at the time of this draft.}

\section{Results}

\subsection{The instrument is reproducible}
\label{sec:repro}

Table~\ref{tab:repeats} gives the same weights run three times. Q8\_0 moved by \qtRepRangeEight{} points across three runs and Q4\_K\_M by \qtRepRangeFour{}. At the level of individual questions, at most \qtRepFlipsMaxEight{} of \qtN{} verdicts changed between two Q8\_0 runs and at most \qtRepFlipsMaxFour{} between two Q4\_K\_M runs. The instrument is close to deterministic. Any difference on the ladder larger than about \qtMaxRepRange{} points is therefore a property of the weights, not of the run.

\begin{table}[t]
\centering\small
\begin{tabular}{@{}llrr@{}}
\toprule
Rung & Accuracy per run & Range (pp) & Max per-question flips \\
\midrule
Q8\_0 & 91.75, 91.50, 91.50 & 0.25 & 1 of 400 \\
Q4\_K\_M & 92.00, 91.50, 92.00 & 0.50 & 4 of 400 \\
\bottomrule
\end{tabular}

\caption{Run-to-run variation with identical weights and settings, \texttt{simple\_python}, $n=\qtN{}$.}
\label{tab:repeats}
\end{table}

\subsection{Flat to four bits, a cliff at three}
\label{sec:ladder}

Table~\ref{tab:ladder} is the ladder. From Q8\_0 down to Q4\_K\_M the accuracy sits in a band of under three points and no rung differs from Q8\_0 at any conventional level: Q4\_K\_M is at \qtAccFour{}\% against Q8\_0's \qtAccEight{}\%, a difference of \qtDeltaFour{} points with a McNemar $p$ of \qtPFour{}, while being \qtSizeRatioFour{}\% of the file size. Q6\_K scores above Q8\_0 and Q5\_K\_M below it, and neither difference survives paired testing, so the apparent non-monotonicity in the raw numbers is within the instrument's noise.

Q3\_K\_M is different in kind. It loses \qtLostThree{} questions that Q8\_0 gets right and recovers \qtGainedThree{}, a net \qtDeltaThree{} points at $p=\qtPThree{}$. That drop is about \qtCliffOverNoise{} times the largest run-to-run range in Table~\ref{tab:repeats}. For this model the tool-calling floor sits between four and three bits.

\begin{table}[t]
\centering\small
\begin{tabular}{@{}lrrrrrr@{}}
\toprule
Rung & Size (GB) & Accuracy & vs.\ Q8\_0 & Lost & Gained & McNemar $p$ \\
\midrule
Q8\_0 & 2.2 & 91.75\% & baseline & & & \\
Q6\_K & 1.7 & 93.50\% & $+$1.75\,pp & 6 & 13 & 0.167 \\
Q5\_K\_M & 1.5 & 90.75\% & $-$1.00\,pp & 10 & 6 & 0.454 \\
Q4\_K\_M & 1.3 & 92.00\% & $+$0.25\,pp & 13 & 14 & 1.000 \\
Q3\_K\_M & 1.1 & 87.00\% & $-$4.75\,pp & 30 & 11 & \textbf{0.004} \\
\bottomrule
\end{tabular}

\caption{The quantization ladder on \texttt{simple\_python}, $n=\qtN{}$ per rung, all rungs from one FP16 source. Lost and gained count questions that flip against Q8\_0; $p$ is McNemar's exact test on those discordant pairs.}
\label{tab:ladder}
\end{table}

\subsection{The free-form floor on the same weights}
\label{sec:freeform}

\blocked{Per-rung GSM8K accuracy under the strict extractor and under the lenient one, the same paired test against Q8\_0, and the location of the free-form floor. Then the verdict on H1: if the free-form floor also sits between Q4\_K\_M and Q3\_K\_M, H1 is not supported and the paper says so; if it sits higher, H1 holds. Also the gap between the two extractors per rung, which bounds how much of the free-form drop is the parser (H3), and median completion length per rung against the chain-of-thought inflation Lotfi et al.\ report.}

\subsection{Across models}
\label{sec:models}

\blocked{The same ladder on the remaining models in the pre-registration: DeepSeek-R1-Distill-Qwen 1.5B and 7B for comparability with Lotfi et al., Qwen3 4B, 8B and 14B, Llama 3.1 8B. The deployment table is this section's product.}

\subsection{Harder categories}

\blocked{Multi-turn, parallel and irrelevance categories on at least the model above. H5 predicts multi-turn degradation at least twice the single-turn relative drop.}

\subsection{Does the overthinking penalty transfer}

\blocked{H4. The training-free logit penalty on overthinking markers from Lotfi et al., applied at Q4\_K\_M and Q3\_K\_M, and its effect on schema validity. Their own scope statement says the marker list is curated for English reasoning models and the evaluation covers only mathematics, code and science, so a null here is informative and is reported either way.}

\section{Limitations}

One model has been measured at the time of this draft, on the easiest \bfcl{} category, on one hardware platform. The ladder is one quantization family; GPTQ, AWQ and the other research methods are deliberately out of scope, and multilingual degradation is excluded because it is already covered. The reproducibility control was run on two rungs rather than all five. \blocked{Anything the free-form and multi-model arms add.}

\section{Conclusion}

\blocked{Written after the free-form and multi-model arms. Shape: for the models measured, the k-quant floor for tool calling sits at [where]; the free-form floor on the same weights sits at [where], so [H1 verdict]; the deployment table gives practitioners the rung to stop at; the one thing a reviewer should take from this is the reproducibility figure, because it is what makes every other number interpretable.}

\section*{Reproducibility}
Every number in this paper expands from \texttt{paper/tables/numbers.tex}, which \texttt{analysis/paper\_numbers.py} recomputes from the raw \bfcl{} result and score files in \texttt{runs/}. The ladder is rebuilt from the FP16 source by \texttt{run\_ladder.sh}, the repeats by \texttt{run\_repeats.sh} and the free-form arm by \texttt{run\_gsm8k.sh}. All raw generations are retained. \blocked{Repository URL and release tag.}

\bibliographystyle{plainnat}
\bibliography{refs}
\end{document}
